\documentclass[12pt,a4paper]{article}

\usepackage[utf8]{inputenc}
\usepackage[T1]{fontenc}
\usepackage[portuguese]{babel}

\usepackage{listings}
\usepackage{xcolor}
\usepackage{amsmath}
\usepackage{amssymb}
\usepackage{amsfonts}
\usepackage{graphicx}

\lstset{
    basicstyle=\ttfamily\small,
    keywordstyle=\color{blue}\bfseries, 
    commentstyle=\color{green!70!black},
    stringstyle=\color{red}, 
    numbers=left, 
    numberstyle=\tiny\color{gray}, 
    frame=single, 
    breaklines=true, 
    showstringspaces=false 
}

\usepackage[margin=2.5cm]{geometry}

\begin{document}

\begin{center}
    {\LARGE MC521 - Relatório II} \\ [2,5em]
    {\large Júlio da Silva Telles} \\ [1,0em]
    {\large 16 de Julho de 2026} \\ [1,5em]
\end{center}

\vspace{1cm} 

\section{G - Vapor Pressure (AtCoder - past202107\_b)}

O problema G, nos da 3 valores distintos $A$, $B$ e $C$ e pede como resposta qual o valor que $A$ deve adquirir tal que,
$C \times B \geq A$, dividido por B.

\subsection{Solução}

A solução desse exercicio pode ser partida em 2 casos:

\begin{itemize}
    \item $A \geq B \times C$: Nesse caso o maior valor que $A$ poderá adquirir será $B \times C$, portanto a resposta será $\frac{B \times C}{B}$ que é apenas $C$.
    \item $A \le B \times C$: Nesse caso a resposta é apenas $A / B$.
\end{itemize}

\section{E - Precontidion (AtCoder - abc412\_b )}

No problema E são dadas duas strings $S$ e $T$, a resposta  para o problema é se a string $S$ cumpre a condição abaixo:

\begin{itemize}
    \item Toda letra maiúscula em $S$, que não esteja no começo da string deve ser immediatamente precedida por um carcter contido na string $T$.
\end{itemize}

\subsection{Solução}

Para resolver esse problema é necessário somente que iteremos sobre todos os indices $i$ de $S$ e a cada vez que encontramos uma letra maiúscula devemos
checar se o carcterer anterior está contido em $T$, para tornar isso mais rápido guardamos todos oc carcteres que aparecem em $T$ em um \texttt{unorderedSet<int>}
que permite a checagem se um elemento existe em $T$ em $\mathcal{O}(1)$.

\end{document}